\documentclass[11pt]{article}
\usepackage{fontspec}
\usepackage[french]{babel}
\usepackage{amsmath}
\usepackage{amssymb}
\usepackage{amsthm}
\usepackage{mathrsfs}
\usepackage{tikz-cd}
\usepackage{enumitem}
\usetikzlibrary{babel}

\newcommand{\printedpage}[1]{\marginpar{\footnotesize p.~#1}}
\newcommand{\msp}[1]{\marginpar{\footnotesize [#1]}}

\newcommand{\Hom}{\mathrm{Hom}}
\newcommand{\Fl}{\mathrm{Fl}}
\newcommand{\Ob}{\mathrm{Ob}}
\newcommand{\Cat}{\mathrm{Cat}}
\newcommand{\Ens}{\mathrm{Ens}}
\newcommand{\pr}{\mathrm{pr}}
\newcommand{\Ker}{\mathrm{Ker}}
\newcommand{\Coker}{\mathrm{Coker}}
\newcommand{\mon}{\mathrm{mon}}
\newcommand{\deq}{\mathrel{\overset{\text{déf}}{=}}}
\newcommand{\id}{\mathrm{id}}
\newcommand{\Winf}{\mathbf{W}_{\infty}}
\newcommand{\dotto}{\mathrel{\tikz[baseline=-0.5ex]\draw[dotted,-latex] (0,0)--(1.6em,0);}}
\begin{document}
% ---- from Der-ChXIIapp-19-06-03_p001-005.tex
\printedpage{}
% no page number printed (title page)
% head: (none)

\begin{center}
{\LARGE\textsf{LES DÉRIVATEURS}}

\bigskip

\textbf{ALEXANDRE GROTHENDIECK}

\bigskip\bigskip

{\large\textbf{Appendice au Chapitre XII}}
\end{center}

\newpage

\printedpage{1}
% head: A. Grothendieck : Dérivateurs, Ch. XII, Édition des Universités de Montpellier II et Paris VII
% foot: Version 20 juin 2003

\msp{page 1}
\textbf{Proposition} (?) (\footnote{Cf. pages 11 ff. pour les contrexemples-clef.}). \textit{Soit $X$ dans} $\Cat$\textit{, avec $b$, $c$ dans $X$ tels que}

\begin{itemize}
\item[\textit{a)}] \textit{$b\times c$ est dans $X$.}
\item[\textit{b)}] \textit{$\forall$ $x$ dans $X$, $b(x)$ ($\deq$} $\Hom_X(x,b)$\textit{) et $c(x)$ sont $\neq\emptyset$, i.e. $X/b \longrightarrow X$ et $X/c \longrightarrow X$ sont surjectifs sur les objets.}
\end{itemize}

\textit{Alors $X$ est $\mathbf{W}_\infty$-asphérique.} (\footnote{Faux, cf. plus bas. $X$ est 1-connexe, mais pas nécessairement 2-connexe.})

Le fait que $b$ soit un plus grand élément de $\Ob X$ assure déjà la 0-connexité.

Pour la 1-connexité, prouvons que dans le groupoïde fondamental $\Pi(X)$, $a \deq b\times c$ devient un élément final.
\[
\begin{tikzcd}
b\times c \arrow[r] \arrow[d] & b \\
c &
\end{tikzcd}
\]

Pour ceci, prouvons d'abord que $\forall$ $x$ dans $X$, tous les éléments de $a(x) = b(x)\times c(x)$ sont égaux dans $\Fl\Pi(X)$. En effet, comme $b\times c \longrightarrow b$ est \texttt{[un]} isomorphisme, donc \texttt{[un]} monomorphisme dans $\Pi(X)$, cela montre déjà que si $\beta,\beta'\in b(x)$, $\gamma,\gamma'\in c(x)$, d'où $(\beta,\gamma)$, $(\beta,\gamma')$, $(\beta',\gamma')\in b(x)\times c(x)$, on a $(\beta,\gamma) = (\beta,\gamma')$ dans $\Pi(X)$ (puisque $\pr_1(\beta,\gamma) = \pr_1(\beta,\gamma')$, et de même (en utilisant que $\pr_2$ est \texttt{[un]} isomorphisme, donc \texttt{[un]} monomorphisme) $(\beta,\gamma') = (\beta',\gamma')$). Donc tous les $(\beta,\gamma)\in b(x)\times c(x) = a(x)$ sont égaux dans $\Fl\Pi(X)$. Pour achever de prouver que $a$ est objet final de $\Pi(X)$, il suffit de voir que toute flèche \textit{dans} $\Pi(X)$ d'un $x\in\Ob X$ dans $a$, provient d'une flèche $x\longrightarrow a$ dans $X$. Mais une flèche dans $\Pi(X)$ est obtenue par un chemin dans $X$
\[
\begin{tikzcd}
x = x_0 \arrow[r, "u_1"] \arrow[dr, "\alpha_0"] & x_1 \arrow[d, "\alpha_1"'] & x_2 \arrow[l, "u_2"'] \arrow[r] \arrow[dl, "\alpha_2"'] & {} & \cdots & {} & x_{2N} = a \arrow[l, "u_{2N}"'] \\
& a. & & & & &
\end{tikzcd}
\]

Prenons une flèche $\alpha_0 : x\longrightarrow a_0$ dans $X$, je dis que la % [sic] a_0 -> a, the arrows α_i go to a (no a_0 is defined)
\msp{page 2}
classe d'homotopie du chemin $c$ (i.e. son image dans $\Fl(\Pi(X))$ est égale à celle de $u_0$. % [sic] u_0 -> \alpha_0 (?), no u_0 is defined; also the parenthesis opened at "(i.e." is not closed
Pour le voir, choisissons pour tout $i\in[0,2N]$ une flèche $\alpha_i : x_i\longrightarrow a$. (Ça existe grâce à l'hypothèse b)). Je dis que pour chaque $i\in[0,2N]$, le triangle
\printedpage{2}
% head: A. Grothendieck : Dérivateurs, Ch. XII, Édition des Universités de Montpellier II et Paris VII
% foot: Version 20 juin 2003
\[
\begin{tikzcd}[column sep=small]
x_0 \arrow[rr, "c_i"] \arrow[dr, "\alpha_0"] & & x_i \arrow[dl, "\alpha_i"'] \\
& a &
\end{tikzcd}
\]
dans $\Pi(X)$ est commutatif, où $c_i$ désigne la classe d'homotopie du chemin de $x_0$ à $x_i$ définie par les $u_i$ ($i\in[1,i]$) \texttt{[plutôt} ($i\in[1,2N]$)\texttt{]}. Nous le faisons par récurrence sur $i$, l'assertion est triviale si $i = 0$ ($c_0$ représente l'application identitique). Le diagramme % [sic] identitique -> identique, non-existent word
\[
\begin{tikzcd}[column sep=small]
x_0 \arrow[rr, "c_i"] \arrow[dr, "\alpha_0"] \arrow[rrrr, bend left=15, "c_{i+1}"] & & x_i \arrow[dl, "\alpha_i"'] \arrow[rr, "u_{i+1}"] & & x_{i+1} \arrow[dlll, "\alpha_{i+1}"'] \\ % ?? uncertain: the curved arrow c_{i+1} -- solid or dashed
& a & & &
\end{tikzcd}
\]
montre qu'il suffit de prouver que pour $i\in[0,2N-1]$ le triangle
\[
\begin{tikzcd}[column sep=small]
x_i \arrow[rr, "u_{i+1}"] \arrow[dr, "\alpha_i"'] & & x_{i+1} \arrow[dl, "\alpha_{i+1}"] \\
& a &
\end{tikzcd}
\]
est commutatif dans $\Pi(X)$. (\textbf{NB} $u_{i+1}$ est en sens directe ou rétrograde, suivant que % [sic] directe -> direct, agreement with "sens"
$i$ est paire ou impaire.) Mais cela résulte du fait que l'image de $\Hom_X(x_i,a)$ (resp. $\Hom_X(x_{i+1},a)$) % [sic] paire ou impaire -> pair ou impair, agreement with "i"
\msp{page 3}
dans $\Fl\Pi(X)$ est réduite à un seul élément.

Ainsi on trouve que $X$ \textit{est} 1-\textit{connexe.} Il y a donc de l'espoir! Il reste à prouver
\[
H^i(X,M) = 0, \qquad i\geq 2, \quad M \text{ coefficient constant.}
\]

J'ai regardé le cas lorsqu'on a ceci : $b = c$, $\Ob X = \{b, b\times b\}$, $\Hom_X(b\times b)$ \texttt{[plutôt} $\Hom(b\times b, b)$\texttt{]} $= \{$ réunion disjointe des ensembles $u\circ\pr_1$ et $u\circ\pr_2$, avec $u\in\Hom_X(b,b) = G$ un monoïde quelconque $\}$. Alors $X^\circ$ est fibrée, via $I\longmapsto b^I$, sur la sous-catégorie pleine $S$ de $\Ens$ formée de $\{0\}$ et $\{0,1\}$, et donc $X$ est cofibrée (en particulier propre) sur $S^\circ \deq T$. (NB $S^\circ = T$ a un objet initial, mais non un objet final. Sa cohomologie à valeurs dans un préfaisceau quelconque n'est pas toujours triviale, cf. ci-dessous.)

On trouve \texttt{[`On trouve' barré]}
\[
H^n(X,M) \Longleftarrow H^p(T, \underbrace{R^q f_*(M_X)}_{= \mathcal{H}^q}), \qquad \text{où } f : X \longrightarrow S^\circ
\]
\texttt{[}$n = p+q$\texttt{]}. J'admet ici (provisoirement) ceci : % [sic] J'admet -> J'admets, verb ending
\[
\begin{aligned}
H^p(T,F) &= 0 \ , &&\text{pour } p\geq 2, \text{ tout } F\in T^\wedge \ ,\\
H^0(T,F) &= \Ker\delta \ , &&\\
H^2(T,F) &= \Coker\delta \ , && % [sic] H^2(T,F) -> H^1(T,F), the first line gives H^p(T,F) = 0 for p >= 2
\end{aligned}
\]
\printedpage{3}
% head: A. Grothendieck : Dérivateurs, Ch. XII, Édition des Universités de Montpellier II et Paris VII
% foot: Version 20 juin 2003
où $\delta : F_0 \longrightarrow F_1$ est défini comme la différence $\delta = p_1^* - p_2^*$, où $p_1$, $p_2$ sont les deux projections de $\{0,1\}\simeq\{0\}^2$ dans $T = S^\circ$. On a de plus
\[
(\mathcal{H}^q)_0 = H^q(G,M), \qquad (\mathcal{H}^q)_1 = H^q(G\times G, M),
\]
et
\[
p_1^*,\ p_2^*\ :\ H^q(G,M) \rightrightarrows H^q(G{\times}G, M) \qquad \text{sont les transposés de } \pr_1,\ \pr_2 : G\times G\longrightarrow G.
\]

\msp{page 4}
On vérifie aisément que
\[
H^0(T,\mathcal{H}^q) = 0 \qquad \text{si } q\geq 1,
\]
et il reste
\[
H^n(X,M) \simeq H^1(S,\mathcal{H}^{n-1}) = \Coker(\delta : H^{n-1}(G) \longrightarrow H^{n-1}(G\times G)). % [sic] H^1(S, ...) -> H^1(T, ...), the cohomology is taken on T = S°
\]
Cela donne, pour $n = 2$,
\[
\begin{aligned}
H^2(X,M) &\simeq \Coker\Big(\delta : \Hom_\mon(G,M) \longrightarrow \underbrace{\mathit{Hom}_\mon(G\times G, M)}_{\Hom_\mon(G,M)\times\Hom_\mon(G,M)} \ \Big)\\
&\simeq \Coker\Big(u\longmapsto (u,-u) : H \longrightarrow H\times H\Big),
\end{aligned}
\]
où $H$ est le groupe abélien $\Hom_\mon(G,M)$.

D'où
\[
H^2(X,M) \ \simeq\ H = \Hom_\mon(G,M).
\]
Donc
\begin{quote}
$X$ est 2-connexe si et seulement si le groupe \textit{abélien} enveloppant du monoïde $G$ est nul.
\end{quote}

Ainsi, si $G$ est un groupe abélien $\neq 0$, $X$ \textit{n'est pas} 2-connexe.

\textbf{Corollaire.} $\Psi = {}$\begin{tikzcd}[ampersand replacement=\&, row sep=tiny, column sep=small, baseline=(A.base)]
\& b_0 \\
|[alias=A]| a_0 \arrow[ur] \arrow[dr] \& \\
\& c_0
\end{tikzcd}
\textit{n'est} pas \textit{totalement} $\mathbf{W}_\infty$\textit{-asphérique; i.e. si on a un diagramme} \begin{tikzcd}[ampersand replacement=\&, row sep=tiny, column sep=small, baseline=(B.base)]
\& b \\
|[alias=B]| a \arrow[ur] \arrow[dr] \& \\
\& c
\end{tikzcd}
\textit{dans une catégorie} $X\in\Ob\Cat$\textit{, tel que sa somme amalgamée dans $X^\wedge$ soit égale à $e_X$, il ne s'ensuit pas que $X$ soit $\mathbf{W}_\infty$-asphérique.} (Mais peut-être voudra bien être 1-connexe?)

\printedpage{4}
% head: A. Grothendieck : Dérivateurs, Ch. XII, Édition des Universités de Montpellier II et Paris VII
% foot: Version 20 juin 2003

\msp{page 5}
\textbf{Question.} \textit{Y-a-t-il d'autres objets totalement} $\mathbf{W}_\infty$\textit{-asphériques de} $\Cat$\textit{, que les catégories ayant un objet final?} Je suspecte à présent que non -- donc que l'on ne peut améliorer substantiellement sur la notion déjà connue d'immersion ouverte de Dwyer.

Dans le diagramme
\[
\begin{tikzcd}
U \arrow[r, hook, "i"] \arrow[d, "p"'] \arrow[dr, phantom, "\substack{\text{cocart.}\\ \text{et cart.}}"] & X \arrow[d, "q"'] & T \arrow[l, hook', dotted, "j"'] \arrow[d, dotted, "\text{iso}"] \\
U' \arrow[r, hook, "i'"] & X' & T', \arrow[l, hook', dotted, "j'"']
\end{tikzcd}
\tag*{$(*)$}
\]
prenons
\[
\begin{aligned}
U = \Psi, \quad & T = \{t\} = e \text{ catégorie finale,}\\
& t \text{ objet final de } X,\\
& \text{i.e. } H_t = (u\longmapsto \Hom_X(u,t)) = e_U \in \Ob U^\wedge \ \mbox{(\footnotemark)}.
\end{aligned}
\]
\footnotetext{i.e. $X$ déduit de $U = \Psi$ en lui rajoutant un objet final $t$ -- donc $X\simeq\boldsymbol{\Delta}^1\times\boldsymbol{\Delta}^1$,
\[
\begin{tikzcd}[ampersand replacement=\&,
execute at end picture={
\draw[rounded corners] ([xshift=-10pt,yshift=10pt]A.center) -- ([xshift=-10pt,yshift=-10pt]C.center) -- ([xshift=10pt,yshift=-10pt]D.center) -- ([xshift=10pt,yshift=3pt]D.center) -- ([xshift=3pt,yshift=3pt]C.center) -- ([xshift=3pt,yshift=10pt]A.center) -- cycle;
\node (P) at ([xshift=-35pt,yshift=-22pt]C.center) {$\Psi$};
\draw[-stealth] (P) -- ([xshift=-11pt,yshift=-4pt]C.center);
}]
|[alias=A]| {} \arrow[r] \& {} \\
|[alias=C]| {} \arrow[u] \arrow[r] \& |[alias=D]| {} \arrow[u]
\end{tikzcd}
\ .
\]}
Prenons pour $U'$ la catégorie dont il est question dans la `proposition' hypothétique (et fausse) p. 1 avec $b = c$, plus précisément une catégorie avec deux objets $b$ et $b\times b$, un groupe (voire un monoïde) $G$ d'automorphismes \texttt{[(voire d'endomorphismes)]} de $b$, et les autres flèches celles qui s'en déduisent, en exigeant que $b\times b$ soit bel et bien le carré cartésien de $b$. On trouve une catégorie cofibrée (donc propre) sur la catégorie $T$ opposée de celle des ensembles finis de la forme $\{0\}$ et $\{0,1\}$ (deux objets) avec comme fibres $B_G$ et $B_{G\times G}$ respectivement ($B_\pi =$ catégorie avec un seul objet, et monoïde d'endomorphismes $\pi$). Le foncteur
\[
U = \Psi \longrightarrow U'
\]
est celui donné par le diagramme

\msp{page 6}
\[
\begin{tikzcd}[row sep=small]
& b \\
b\times b \arrow[ur, "\pr_1"] \arrow[dr, "\pr_2"'] & \\
& b
\end{tikzcd}
\]

% ---- from Der-ChXIIapp-19-06-03_p006-010.tex
\printedpage{5}
% head: A. Grothendieck : Dérivateurs, Ch. XII, Édition des Universités de Montpellier II et Paris VII
% foot: Version 20 juin 2003

dans $U'$, et les hypothèses a), b) de la page 1 assurent que la somme amalgamée $\varinjlim$ de ce diagramme dans $U'^{\wedge}$ est $e_{U'}$, i.e. que
\[ p_!(e_U) = e_{U'}. \]
Il s'ensuit que dans la somme amalgamée, $X'$ envisagée dans $(*)$, l'unique élément $t'$ de $T'$ est objet final de $X'$. On a vu que $U'$ est 1-connexe, mais non 2-connexe si le groupe abélien enveloppant de $G$ (i.e. $G/$(groupe des commutateurs) dans le cas où $G$ est un groupe) n'est pas réduit à 1. L'exemple le plus économique s'obtient en prenant $G = \{\pm 1\}$ le groupe à deux éléments. C'est donc un cas où dans $(*)$ on a
\[ i \in W \ \ \text{(et même $i$ asphérique),} \qquad \text{mais } i' \notin W, \]
et qui a l'avantage sur l'exemple plus sophistiqué de la page 103 (prop. 10) de ne mettre en jeu que des catégories finies -- et même fort petites !

Procédons maintenant comme dans la proposition 11, en désignant par $U''$ la catégorie déduite de $U'$ en lui ajoutant un objet
\msp{page 7}
final, et soit
\[ p' : U' \longrightarrow U'' \]
l'inclusion, de sorte que
\[
\begin{aligned}
p'_!(e_{U'}) \quad &= \quad (p'p)_!(e_U) = U' \quad && \text{en tant qu'objet dans } U''^{\wedge}\\
& && \text{(sous-préfaisceau du préfaisceau final } e_{U''}\text{).}
\end{aligned}
\]
Considérons
\[
\begin{tikzcd}
U \arrow[r, hook, "i"] \arrow[d, "{p''=p'p}"'] \arrow[dr, phantom, "\begin{smallmatrix}\text{cocart.}\\ \text{et cart.}\end{smallmatrix}"] & X \arrow[d, "q''"'] & T \arrow[l, hook', dotted, "j"'] \arrow[d, dotted, "\text{iso}"] \\
U'' \arrow[r, hook, "i''"] & X'' & T''. \arrow[l, hook', dotted, "j''"']
\end{tikzcd}
\]
Cette fois $U$, $X$ \textit{et} $U''$ sont $\Winf$-asphériques (et même contractiles toutes trois), donc $i$, $p''$ sont dans $W$, pourtant $X''$ n'est pas $\Winf$-asphérique (donc $i'', q'' \notin \Winf$), car pour qu'elle le soit, i.e. pour que $i'' \in \Winf$, il faudrait (prop. 7, cor. page 85 (\footnote[4]{pas encore démontré !})) que $i''$ soit $\Winf$-asphérique, i.e. que $U''/t''$ soit $\Winf$-asphérique. Mais par construction $U''/t'' \simeq U'$, qui n'est \textit{pas} $\Winf$-asphérique.

Voilà un exemple encore plus simple d'un foncteur $f : \Psi \longrightarrow X$ tel que $f_!(e_{\Psi})$ ne soit pas $\Winf$-asphérique, i.e. d'un diagramme
\[
\begin{tikzcd}
& b \\
a \arrow[ur, "\beta_0"] \arrow[dr, "\gamma_0"'] & \\
& c
\end{tikzcd}
\]
\printedpage{6}
% head: A. Grothendieck : Dérivateurs, Ch. XII, Édition des Universités de Montpellier II et Paris VII
% foot: Version 20 juin 2003
\msp{page 8}
dans une catégorie $V$, telle que la limite inductive $L = b \amalg_a c$ \textit{dans} $V^{\wedge}$ soit un objet non asphérique de $V^{\wedge}$, i.e. $V/L$ soit non asphérique. Je sais qu'il faut, pour un tel contre-exemple, que $a \longrightarrow b$ et $a \longrightarrow c$ ne soient pas des monomorphismes. L'idée, c'est de prendre un exemple universel où ces flèches ne soient pas \texttt{[des]} monomorphismes, en prenant
\msp{page 9}
\[
V = \left(
\begin{tikzcd}
& & b \\
0 \arrow[r, shift left, "\lambda"] \arrow[r, shift right, "\mu"'] & a \arrow[ur, "\beta"] \arrow[dr, "\gamma"'] & \\
& & c
\end{tikzcd}
\right)
\]
avec
\[ \beta\lambda = \beta\mu \qquad \text{et} \qquad \gamma\lambda = \gamma\mu \]
les deux seules flèches non identiques, en plus des quatre flèches marquées $\beta$, $\gamma$, $\lambda$, $\mu$.

On trouve que
\[ b \amalg_a c \xrightarrow{\ \sim\ } e_X \qquad \text{(somme amalgamée dans } V^{\wedge}\text{).} \] % [sic] e_X -> e_V, the category considered here is V

Pourtant $V$ n'est pas asphérique, ni même 1-connexe : il est homotope à la sous-catégorie ouverte $V_0 = (0 \overset{\lambda}{\underset{\mu}{\rightrightarrows}} a)$, car on a une rétraction $V \xrightarrow{\ r\ } V_0$ (envoyant $b$ et $c$ sur $a$), et une flèche $ir \longrightarrow \id_V$ (où $i : X_0 \hookrightarrow X$ est l'inclusion). Donc $V$ a le même type d'homotopie \textit{que} $V_0$, qui est un cercle. Donc on a intérêt à prendre ce $V$, où $\tilde{V}$ déduit de $V$ en rajoutant un élément final $e$, % [sic] i : X_0 \hookrightarrow X -> i : V_0 \hookrightarrow V, the inclusion of V_0 in V is meant
\[
\tilde{V} = \left(
\begin{tikzcd}
& & b \arrow[dr] & \\
0 \arrow[r, shift left] \arrow[r, shift right] & a \arrow[ur] \arrow[dr] & & e \\
& & c \arrow[ur] &
\end{tikzcd}
\right),
\]
comme $U'$, pour les contre-exemples pages 6 et 7.

\msp{page 10}
\subsection*{2.6. Contre-exemples pour la totale asphéricité.}

Je me suis finalement convaincu que
\[
\Psi = \left\{
\begin{tikzcd}
& b_0 \\
a_0 \arrow[ur] \arrow[dr] & \\
& c_0
\end{tikzcd}
\right\}
\]
\printedpage{7}
% head: A. Grothendieck : Dérivateurs, Ch. XII, Édition des Universités de Montpellier II et Paris VII
% foot: Version 20 juin 2003
n'est \textit{pas} totalement asphérique (pour le localiseur $\Winf$ -- dans le présent numéro, on supposera $W = \Winf$), ou ce qui revient au même, qu'on peut avoir un diagramme de type $\Psi$
\[
\begin{tikzcd}
& b \\
a \arrow[ur, "\beta"] \arrow[dr, "\gamma"'] & \\
& c
\end{tikzcd}
\tag*{$(*)$}
\]
dans une catégorie $Z \in \mathrm{Ob}\,\mathrm{Cat}$, dont la limite inductive $b \amalg_a c$ dans $Z^{\wedge}$ soit le préfaisceau final $e_Z$, i.e. tel que l'on ait
\[ b(z) \amalg_{a(z)} c(z) \xrightarrow{\ \sim\ } \text{ensemble ponctuel} \]
pour tout objet $z$ de $Z$ (où $a(z) = \Hom_Z(a,z)$ \texttt{[plutôt} $\Hom_Z(z,a)$\texttt{]}) sans que pour autant $Z$ soit asphérique, ni même 1-connexe (\footnote[5]{faux -- il est 1-connexe forcément (\footnotemark[6]), mais non 2-connexe.}).\footnotetext[6]{cf. commentaire XII, p. 101.} L'exemple le plus simple que j'ai trouvé, est l'exemple ‘universel’ d'un diagramme $(*)$ où ni l'une ni l'autre des deux flèches $a \longrightarrow b$, $a \longrightarrow c$ n'est \texttt{[un]} monomorphisme (car si elle l'est, alors on peut prouver que $Z$ est bel et bien asphérique, comme nous verrons par la suite). C'est la catégorie
\[
Z = \left\{
\begin{tikzcd}
& & b_0 \\
0_0 \arrow[r, shift left, "\lambda_0"] \arrow[r, shift right, "\mu_0"'] & a_0 \arrow[ur, "\beta_0"] \arrow[dr, "\gamma_0"'] & \\
& & c_0
\end{tikzcd}
\right\}
\]
\msp{page 11}
déduite de $\Psi$ en lui rajoutant un objet $0_0$, relié à $a_0$ par deux flèches distinctes $\lambda_0$, $\mu_0$, telles que
\[ \beta_0\lambda_0 = \beta_0\mu_0 \qquad \text{et} \qquad \gamma_0\lambda_0 = \gamma_0\mu_0, \]
ces deux flèches composées étant, avec les quatre flèches $\beta_0$, $\gamma_0$, $\lambda_0$, $\mu_0$ marquées sur le diagramme, les seules (six) flèches non identiques de $Z$. $\Psi$ immergée dans $Z$ comme un fermé, de complémentaire la catégorie ponctuelle réduite à $0_0$. (NB L'immersion fermée $\Psi \longrightarrow Z$ n'est pas immersion de Dwyer, i.e. (comme $Z$ est le plus petit ouvert de $Z$ contenant $\Psi$) n'est pas coasphérique, car la catégorie $0\backslash\Psi$, cofibrée sur $\Psi$ avec fibres discrètes, de cardinal 2 en $a_0$ et 1 en $b_0$ et $c_0$, est la catégorie
\[
\begin{tikzcd}
& \tilde{b}_0 \\
a'_0 \arrow[ur] \arrow[ddr] & \\
a''_0 \arrow[uur, crossing over] \arrow[dr] & \\
& \tilde{c}_0
\end{tikzcd}
\]
\printedpage{8}
% head: A. Grothendieck : Dérivateurs, Ch. XII, Édition des Universités de Montpellier II et Paris VII
% foot: Version 20 juin 2003
i.e.
\[
\begin{tikzcd}
& {} & \\
{} \arrow[ur] \arrow[dr] & & {} \arrow[ul] \arrow[dl] \\
& {} &
\end{tikzcd}
\quad ,
\]
qui a le type d'homotopie du cercle.) Ni l'immersion ouverte $\{0_0\} \longrightarrow Z$, car $\{0_0\}/a_0$ est la catégorie discrète de cardinalité 2.

L'inclusion $i : Z_0 \longrightarrow Z$ de
\[ Z_0 \ = \ (\ 0_0 \overset{\lambda_0}{\underset{\mu_0}{\rightrightarrows}} a_0\ ) \]
dans $Z$ est un homotopisme (\footnote[7]{faux, il n'y a pas de rétraction du tout! En fait, $Z$ est la \textit{suspension} de la $\rightrightarrows$, qui a le type d'homotopie (topologique) du cercle, et il a donc le type d'homotopie de la 2-sphère.}) (car il existe une rétraction évidente $r$ de $Z$ sur $Z_0$ (\footnotemark[7]), et une flèche $ir \longrightarrow \id_Z$), donc $Z$ est homotope à $Z_0$, qui a le type d'homotopie du cercle, donc n'est pas 1-connexe.

\msp{page 12}
D'autre part, la vérification de
\[ b_0 \amalg_{a_0} c_0 \ = \ e_Z \qquad \text{dans } Z^{\wedge} \]
est immédiate.

On va utiliser cette flèche
\[ \Psi \longrightarrow Z \]
pour construire des exemples intéressants (et déchirants) de sommes amalgamées
\[
\begin{tikzcd}
U \arrow[r, hook, "i"] \arrow[d, "p"'] & X \arrow[d, "q"'] \\
U' \arrow[r, hook, "i'"] & X',
\end{tikzcd}
\]
avec $i$ immersion ouverte, telle qu'on ait même
\[ i :\ U \longrightarrow \overline{U} \qquad \text{est asphérique, donc dans } W, \]
alors que pourtant on aura
\[ i' :\ U' \longrightarrow \overline{U'} \qquad \notin W, \]
et ce qui mieux (ou pis) est,
\[ p \in W \qquad \text{mais} \qquad q \notin W \]
\printedpage{9}
% head: A. Grothendieck : Dérivateurs, Ch. XII, Édition des Universités de Montpellier II et Paris VII
% foot: Version 20 juin 2003
(et éventuellement, en même temps, $i \in W$ mais $i' \notin W$). De plus (contrairement aux constructions des propositions 10, 11 (p. 101-107)), les catégories envisagées pourront toutes être finies. Je suis particulièrement intéressé par le cas où $X$ est un ‘carré’
\[ X = \Phi_n \times \Phi_m \qquad (n,m \geq 2), \]
et où $U$ est son ‘bord’
\[ U = \dot{X} \overset{\text{déf}}{=} (\dot{\Phi}_n \times \Phi_m) \cup (\Phi_n \times \dot{\Phi}_m). \]

\msp{page 13}
\[
\begin{tikzcd}
0 \arrow[r] \arrow[d] & a \arrow[d] & a' \arrow[l] \arrow[r] \arrow[d] & \cdots \\
b \arrow[r] & t_1 & \cdots \arrow[l] & \\
b' \arrow[u] \arrow[r] \arrow[d] & {} \arrow[u] & \cdots \arrow[l] \arrow[u] & \\
\cdots & & &
\end{tikzcd}
\]

\textbf{Proposition.} \textit{Soient $n,m \geq 2$, et considérons l'inclusion $U = \partial X \overset{i}{\hookrightarrow} X$ du bord du ‘carré’ $X = \Phi_n \times \Phi_m$ dans $X$. C'est une immersion ouverte $W$-asphérique, de façon précise, les $U/t$ ($t \in \dot{U} = \bar{U}\backslash U$) sont isomorphes à $e$ (catégorie ponctuelle) ou à $\Phi =$ \begin{tikzcd}[row sep=small, column sep=small] & {} \arrow[dl] \\ {} & \\ & {} \arrow[ul] \end{tikzcd} (catégorie avec élément final), sauf pour les quatre éléments $t_i$ des quatre ‘coins’ de $\dot{U}$, où c'est la catégorie}
\[
\tilde{\Psi} = \left(
\begin{tikzcd}
& & a' \arrow[dl] \\
& a & \\
0 \arrow[ur] \arrow[dr] & & \\
& b & \\
& & b' \arrow[ul]
\end{tikzcd}
\right)
\]
\textit{n'ayant pas d'élément final. Donc $U \xrightarrow{\ i\ } X$ n'est pas immersion ouverte de Dwyer (mais $U \longrightarrow X\backslash\{t_1,t_2,t_3,t_4\}$ l'est \ldots). Ceci posé, je dis que $i : U \longrightarrow X$ n'est pas une $W$-cofibration.}

% ---- from Der-ChXIIapp-19-06-03_p011-014.tex
\printedpage{10}
% head: A. Grothendieck : Dérivateurs, Ch. XII, Édition des Universités de Montpellier II et Paris VII
% foot: Version 20 juin 2003
\textsc{Démonstration}. Soit
\[ F \hookrightarrow U \]
l'inclusion dans $U = \partial X$ de la catégorie-coin \texttt{[}$V =$\texttt{]} $\vcenter{\hbox{\begin{tikzcd}[row sep=small, column sep=small]
0 \arrow[r] \arrow[d] & a \\
b &
\end{tikzcd}}}$ (cf. figure), isomorphe à $\Psi$. On voit aussitôt que c'est une immersion fermée de Dwyer (le plus petit ouvert de $U$ contenant $\Psi$ est justement la catégorie $\tilde{\Psi} = U/t_1$ de tantôt), donc une $W$-cofibration. Si $U \longrightarrow X$ était une $W$-cofibration, il en serait de même du composé $V \longrightarrow X$. On va voir que ce n'est pas le cas.

Soit $\overline{V} = \vcenter{\hbox{\begin{tikzcd}[row sep=small, column sep=small]
0 \arrow[r] \arrow[d] & a \arrow[d] \\
b \arrow[r] & t_1
\end{tikzcd}}}$ l'adhérence de $V$ dans $X$,
\msp{page 14}et considérons l'inclusion
\[ \overline{V} \longrightarrow X, \]
je dis que \texttt{[c]}'est une immersion fermée de Dwyer, i.e. que pour tout $s \in W\backslash\overline{V}$, où $W$ est le plus petit ouvert de $X$ contenant $\overline{V}$, i.e.
\[
W = \quad
\begin{tikzcd}[row sep=small, column sep=small]
0 \arrow[r] \arrow[d] & a \arrow[d] & a' \arrow[l] \arrow[d] \\
b \arrow[r] & {} & a'' \arrow[l] \\
b' \arrow[u] \arrow[r] & b'' \arrow[u] & c, \arrow[l] \arrow[u]
\end{tikzcd}
\]
la catégorie $s\backslash\overline{V}$ a un objet initial. (Les catégories $s\backslash\overline{V}$ qu'on obtient, à partir des cinq choix possibles pour $s$, sont $\boldsymbol{\Delta}^1$ (si $s = a'$ ou $b'$) ou $e$ (pour $s = a''$, $b''$ ou $c$)). Donc $\overline{V} \longrightarrow X$ est une $W$-cofibration.

Ceci posé, on sait qu'il existe, pour [?] l'inclusion (immersion ouverte)
\[ V \hookrightarrow \overline{V} \qquad \text{i.e. } \Psi \hookrightarrow \boldsymbol{\Delta}^1 \times \boldsymbol{\Delta}^1, \]
un foncteur $V \longrightarrow V'$, donnant lieu à un carré cocartésien
\[
\begin{tikzcd}
V \arrow[r, hook, "i"] \arrow[d, "p"'] & \overline{V} \arrow[d, "\overline{p}"] \\
V' \arrow[r, hook, "i'"] & \overline{V'}
\end{tikzcd}
\tag*{$(*)$}
\]
avec
\[ i, p \in W \ \text{(car $V$, $V'$, $\overline{V}$ $W$-asphérique), mais } i', \overline{p} \notin W \ \text{(i.e. $\overline{V'}$ \textit{pas} $W$-asphérique).} \]

Je rappelle qu'on prend d'abord
\printedpage{11}
% head: A. Grothendieck : Dérivateurs, Ch. XII, Édition des Universités de Montpellier II et Paris VII
% foot: Version 20 juin 2003
\[
V'_0 = \left(
\begin{tikzcd}[row sep=small]
 & & a \\
0' \arrow[r, shift left, "\lambda"] \arrow[r, shift right, "\mu"'] & 0 \arrow[ur, "\alpha"] \arrow[dr, "\beta"'] & \\
 & & b
\end{tikzcd}
\right)
\qquad \text{avec} \quad
\left\{
\begin{aligned}
\alpha\lambda &= \alpha\mu \\
\beta\lambda &= \beta\mu
\end{aligned}
\right.
\]
et l'immersion fermée ‘identique’ de $V$ dans $V'$ \texttt{[plutôt} $V'_0$\texttt{]}
\[ p_0 : V \longrightarrow V'_0 \]
\msp{page 15}(qui est telle que
\[ p_{0!}(e_V) = e_{V'_0} \ \ ), \]
puis on prend pour $V'$ la catégorie obtenue en ajoutant un objet final à $V'_0$
\[
V' = \left(
\begin{tikzcd}[row sep=small, column sep=small]
 & & a \arrow[dr, "\rho"] & \\
0' \arrow[r, shift left, "\lambda"] \arrow[r, shift right, "\mu"'] & 0 \arrow[ur, "\alpha"] \arrow[dr, "\beta"'] & & e \\
 & & b \arrow[ur, "\sigma"'] &
\end{tikzcd}
\right)
\qquad \text{avec} \quad
\left\{
\begin{aligned}
\alpha\lambda &= \alpha\mu \\
\beta\lambda &= \beta\mu \\
\rho\alpha &= \sigma\beta
\end{aligned}
\right.
\]
et pour $p$ le foncteur composé
\[
\begin{tikzcd}
V \arrow[r, "p_0"] \arrow[rr, bend right, "p"'] & V'_0 \arrow[r, "p_1"] & V',
\end{tikzcd}
\]
où $p_1$ est l'inclusion (ouverte). On trouve alors le carré $(*)$, avec $\overline{V'}$ non $W$-asphérique, ou ce qui revient au même, $i' \notin W$. Car comme $\overline{V'}\backslash V' \simeq \overline{V}\backslash V$ est la catégorie ponctuelle, $i' \in W$ équivaut à $V'/t$ $W$-asphérique (\footnote[8]{moyennant (W 7 bis).}) (cf. cor. 4, p. 149 dans XII). Or par construction $V'/t \simeq V'_0$, qui n'est pas $W$-asphérique (dans le cas, disons, où $W \subseteq W_1$, car $\pi_1(V'_0) \simeq \mathbf{Z}$).

\msp{page 16}Considérons alors le diagramme de carrés cocartésiens
\[
\begin{tikzcd}
V \arrow[r, hook, "i"] \arrow[d, "p"'] & \overline{V} \arrow[r, hook] \arrow[d, "\overline{p}"] & X \arrow[d, "q"] \\
V' \arrow[r, hook, "i'"] & \overline{V'} \arrow[r, hook] & X'
\end{tikzcd}
\qquad .
\]

\printedpage{12}
% head: A. Grothendieck : Dérivateurs, Ch. XII, Édition des Universités de Montpellier II et Paris VII
% foot: Version 20 juin 2003
Comme $\overline{V} \hookrightarrow X$ est \texttt{[une]} immersion fermée de Dwyer, le carré de droite est $W$-cocartésien, donc $(\overline{p} \in W) \iff (q \in W)$, donc (comme $\overline{p} \notin W$) on a
\[ p \in W, \qquad \text{mais} \qquad q \notin W, \]
ce qui montre bien que $V \longrightarrow X$ n'est pas une $W$-cofibration. Factorisant $V \longrightarrow X$ par $U$ (au lieu de $\overline{V}$), on trouve
\[
\begin{tikzcd}
V \arrow[r] \arrow[d, "p\in W"'] & U \arrow[r] \arrow[d, "q_U\in W"] & X \arrow[d, "q\notin W"] \\
V' \arrow[r] & U' \arrow[r] & X'
\end{tikzcd}
\qquad ,
\]
où $q_U \in W$ puisque $p \in W$ et $V \longrightarrow U$ est \texttt{[une]} immersion de Dwyer. Cela rend explicite que $U \longrightarrow X$ n'est pas une $W$-cofibration, puisque
\[ q_U \in W, \qquad \text{mais} \qquad q \notin W, \]
\hfill q.e.d.

\msp{page 17}\textbf{Recollement dans} Cat \textbf{d'un ouvert et d'un fermé complémentaire.}

\texttt{[Supposons donné]}
\[ X_0 \underbrace{\hookrightarrow}_{\text{ouvert}} X \ \underbrace{\hookleftarrow}_{\substack{\text{fermé}\\ \text{complé-}\\ \text{mentaire}}} X_1 \]
dans Cat, \texttt{[et soit]}
\[ \tilde{X}_{01} \subseteq \mathrm{Fl}(X) \]
\texttt{[la]} sous-catégorie pleine des flèches de source dans $X_0$, but dans $X_1$.

\textbf{Question.} \textit{Ce diagramme}
\[
\begin{tikzcd}[column sep=small]
 & & \tilde{X}_{01} \arrow[dll, "p"' near start, "\text{fibrant, fibres les $x_0\backslash X_1$}"' near end] \arrow[drr, "q" near start, "\text{cofibrant, fibres les $X_0/x_1$}" near end] & & \\
X_0 \arrow[ddrr, "i"'] & & & & X_1 \arrow[ddll, "j"] \\
 & {} \arrow[rr, bend left] & & {} & \\
 & & X, & &
\end{tikzcd}
\tag*{$(*)$}
\]
\printedpage{13}
% head: A. Grothendieck : Dérivateurs, Ch. XII, Édition des Universités de Montpellier II et Paris VII
% foot: Version 20 juin 2003
\textit{est-il $W$-cocartésien?}

OK si $i$ $W$-asphérique ou j $W$-coasphérique, i.e. $X \longrightarrow \boldsymbol{\Delta}^1$ lisse ou propre -- c'est un cas où on a soit $X_0 \xrightarrow{W} X$, soit $X_1 \xrightarrow{W} X$. Mais même en supposant $i$ ou $j$ dans $W$, et de plus $X$ asphérique (donc $X_0$ et $X$ asphériques, ou $X_1$ et $X$ asphériques), et $W = \mathbf{W}_\infty$, je n'ai pu prouver que le diagramme est cocartésien, i.e. $q \in W$ resp. $p \in W$ (\footnote[9]{Je suppose que $W$ satisfait Loc (6).}). (Pas non plus si $X_0 \simeq e$, $X$, $X_1$ $W$-asphérique $\Longrightarrow$ $e\backslash X$ $W$-asphérique, i.e. $X_1 \longrightarrow X$ $W$-coasphérique $\ldots$\,) Mais j'ai prouvé du moins, dans tous les cas, que % [sic] j (upright) -> $j$, letter denoting a functor printed in roman type
\[
\begin{tikzcd}[row sep=small]
 & X_0 & & \\
\mathcal{X} \overset{\text{déf}}{=} \displaystyle\int \tilde{X}_{01} \arrow[ur] \arrow[dr] & & {} \arrow[r, "\varphi"] & X \\
 & X_1 & &
\end{tikzcd}
\tag*{$(**)$} % ?? uncertain: script letter X -- \mathcal{X} or \mathscr{X}
\]
est dans $\mathbf{W}_1$.

Chose \textit{fausse} (sauf si $i$ ou $j$ dans $W = \mathbf{W}_\infty$) si on remplace $\tilde{X}_{01}$ par $\tilde{X}'_{01}$, formée des triples $(x_0, x_1, c)$, avec $x_0 \in \mathrm{Ob}\,X_0$, $x_1 \in \mathrm{Ob}\,X_1$, $c : x_0 \dotto x_1$ un \textit{chemin} de $x_0$ à $x_1$ dans $X$! On a
\msp{page 18}donc une situation ‘topologique’ tout à fait spéciale à la catégorie de modèles Cat, qui ne semble pas avoir d'analogue dans des catégories telles que Esp ou $\boldsymbol{\Delta}^\wedge$, où pourtant les notions d'ouvert et de fermé complémentaire (avec grain de sel dans $\boldsymbol{\Delta}^\wedge$) ont un sens.

L'argument catégorique quasi-trivial qui établit que $(**)$ est dans $\mathbf{W}_1$, i.e. induit une équivalence de catégories entre $\mathrm{Rev}(X)$ et
\[ \mathrm{Rev}(\mathcal{X}) \ \simeq \ \mathrm{Rev}(X_0) \overset{2}{\times}_{\mathrm{Rev}(\tilde{X}_{01})} \mathrm{Rev}(X_1), \]
me semble devoir s'étendre aux dimensions supérieures, à coups de $n$-gerbes sur $X$ -- mais les fondements sur la théorie des $n$-gerbes manquent. (Et je n'ai pas eu le courage d'expliciter l'argument même pour $n = 2$ (où pourtant \textsc{Giraud} fournit ce qu'il faut) (\footnote[10]{Mais j'ai fait quand même une esquisse de vérification pour $n = 2$, qui me convainc qu'on a bien $\varphi \in \mathbf{W}_2$. Et je n'ai pas guère de doute qu'on a même $\varphi \in \mathbf{W}_\infty$.}).) Aussi je conjecture que $\varphi$ est tout au moins dans $\mathbf{W}_\infty$, i.e. que le diagramme $(*)$ est $W$-cocartésien si $W \supseteq \mathbf{W}_\infty$ (et satisfait Loc (6 a$'$)). % [sic] je n'ai pas guère (footnote 10) -> je n'ai guère, double negation pas + guère
\end{document}
